% Part: computability
% Chapter: recursive-functions
% Section: pr-functions-computable

\documentclass[../../../include/open-logic-section]{subfiles}

\begin{document}

\olfileid[tr]{cmp}{rec}{cmp}
\olsection{İlkel Özyinelemeli Fonksiyonlar Hesaplanabilirdir}

Bir $h$ fonksiyonu ilkel özyinelemeyle
\begin{eqnarray*}
  h(\vec x,0) & = & f(\vec x) \\
  h(\vec x,y+1) & = & g(\vec x,y,h(\vec x,y))
\end{eqnarray*}
biçiminde tanımlansın ve $f$ ile $g$ fonksiyonları hesaplanabilir olsun.
($x_0,\dots,x_{k-1}$ dizisini kısaltmak için $\vec x$ yazıyoruz.)
$h(\vec x,0)$ yalnızca hesaplanabilir olduğunu varsaydığımız $f(\vec x)$
olduğundan açıkça hesaplanabilir. Ardından $h(\vec x,1)$ de hesaplanabilir;
çünkü $1=0+1$ ve
\begin{align*}
  h(\vec x,1) & = g(\vec x,0,h(\vec x,0))=g(\vec x,0,f(\vec x)).
  \intertext{Bu biçimde sürdürerek şunları hesaplayabiliriz:}
  h(\vec x,2) & = g(\vec x,1,h(\vec x,1))
    =g(\vec x,1,g(\vec x,0,f(\vec x))) \\
  h(\vec x,3) & = g(\vec x,2,h(\vec x,2))
    =g(\vec x,2,g(\vec x,1,g(\vec x,0,f(\vec x)))) \\
  h(\vec x,4) & = g(\vec x,3,h(\vec x,3))
    =g(\vec x,3,g(\vec x,2,g(\vec x,1,g(\vec x,0,f(\vec x))))) \\
  & \vdots
\end{align*}
Dolayısıyla genel olarak $h(\vec x,y)$ değerini hesaplamak için
$h(\vec x,0),h(\vec x,1),\dots$ değerlerini art arda hesaplar ve
$h(\vec x,y)$ değerine ulaşınca dururuz.

Demek ki $f$ ile $g$ hesaplanabilirse ilkel özyinelemeli bir tanım yeni bir
hesaplanabilir fonksiyon verir. Fonksiyonların bileşkesi de $f$ ile $g_i$
fonksiyonları hesaplanabilir olduğunda hesaplanabilir bir fonksiyon verir.

Temel fonksiyonlar $\Zero$, $\Succ$ ve $\Proj{n}{i}$ hesaplanabilir ve
bileşke ile ilkel özyineleme hesaplanabilir fonksiyonlardan hesaplanabilir
fonksiyonlar ürettiği için her ilkel özyinelemeli fonksiyon hesaplanabilirdir.

\end{document}

